\documentclass[11pt]{article}
\usepackage[margin=1in]{geometry}
\usepackage{amsmath,amssymb,booktabs}
\usepackage[hidelinks]{hyperref}
\newcommand{\Caputo}[1]{{}^{C}_{0}D_t^{#1}}
\title{PDE-RSI Benchmark Problems}
\author{Generated by \texttt{scripts/write\_benchmarks\_latex.py}}
\date{\today}
\begin{document}
\maketitle
\tableofcontents

\section{Protocols, data and scoring}

\paragraph{Student protocol.} For the \texttt{pde\_benchmarks.tex} suites and the inventory suite, an agent receives,
for each case, 256 noisy samples of the solution (Gaussian noise with standard deviation 1\,\% of the samples' RMS,
per component), the domain, and black-box forcing, initial, boundary and (where present) coefficient data. It must
identify the equation class and its coefficients from a library of candidate equations and return the solution over
the whole domain. Case identifiers are hash-ordered and reveal neither the class nor the problem.

\paragraph{Code protocol.} For the manufactured suites (Section~\ref{sec:manufactured}) the equation is given and the
agent writes one solver per family.

\paragraph{Scoring.} The returned solution is linearly interpolated onto a uniform evaluation grid and compared with
the exact or a high-accuracy reference solution using the trapezoidal rule. A case passes when the space--time and
the final-time relative $L^2$ errors are both at most the tolerance, the initial and boundary conditions hold to
$10^{-10}$, and the solve fits the time limit (wall and CPU time, including start-up) and the memory limit. The
reported class and coefficients are scored separately and never gate a case.

\paragraph{Splits.} Training (evolution) cases come with labelled feedback; validation (development) cases drive the
acceptance gate; sealed test cases are never shown to any agent and are scored once, after the solver is frozen.
\begin{center}
\small
\begin{tabular}{p{3.3cm}p{2.3cm}p{2.1cm}p{2.4cm}p{3.6cm}}
\toprule
Suite & Cases per split & Tolerance & Limits & Evaluation grid \\
\midrule
\texttt{pde\_benchmarks.tex} 1--11 & 16 / 16 / 17, 20, 20 & $2\times10^{-3}$ & 90 s, 16 GB & 1-D $512\times500$; 2-D $128^2\times100$ \\
\texttt{pde\_benchmarks.tex} 12--16 & 10 / 10 / 10 & $2\times10^{-3}$ & 90 s (1-D), 300 s (2-D), 16 GB & 1-D $512\times501$; 2-D $129^2\times21$ \\
PDE Data Inventory & 12 / 12 / 12, 12 & $2\times10^{-3}$ (4-D: $5\times10^{-2}$) & 90 s / 300 s, 16 GB & 1-D $512\times500$; 2-D $256^2$ or $128^2\times11$ \\
Manufactured (code protocol) & 24 / 24 / 24 (1-D) & $2\times10^{-3}$ & 15 s, 1 GB & 1-D $512\times500$; wave $128^2\times100$ \\
\bottomrule
\end{tabular}
\end{center}
Vector solutions (Problems 13, 15, 16) are scored with all components together; the pressure is compared after removing its spatial mean.
\section{\texttt{pde\_benchmarks.tex}: Problems 1--16}
Problems 1--11 form the \texttt{pdebench} suite and Problems 12--16 the \texttt{pdebench-ext} suite. The definitions below are those of \texttt{pde\_benchmarks.tex}.
\setcounter{subsection}{0}
\renewcommand{\thesubsection}{\arabic{subsection}}
\subsection{Preliminaries}
%=====================================================================

\subsubsection*{Spatial Profiles}

Cases 1a--4a and Problems 5, 7, and 8 are posed on $x\in(0,L)$ (on
$(0,L)^2$ for Problem 5) with homogeneous Dirichlet boundary conditions.
The remaining cases of Problems 1--4 specify their own domains and
boundary data. The spatial
profile $s$ is one of the following:
\[
s(x)=
\begin{cases}
\sin(\pi x/L),
& \text{training},\\[4pt]
4\,\dfrac{x}{L}\left(1-\dfrac{x}{L}\right),
& \text{validation},\\[6pt]
\sin(\pi x/L)+\dfrac14\sin(3\pi x/L),
& \text{reserved}.
\end{cases}
\]
Each profile satisfies $s(0)=s(L)=0$. The derivatives are
\[
\begin{array}{lll}
\text{training:} &
s'=\dfrac{\pi}{L}\cos\dfrac{\pi x}{L}, &
s''=-\dfrac{\pi^2}{L^2}\sin\dfrac{\pi x}{L},\\[10pt]
\text{validation:} &
s'=\dfrac{4}{L}\left(1-\dfrac{2x}{L}\right), &
s''=-\dfrac{8}{L^2},\\[10pt]
\text{reserved:} &
s'=\dfrac{\pi}{L}\cos\dfrac{\pi x}{L}+\dfrac{3\pi}{4L}\cos\dfrac{3\pi x}{L}, &
s''=-\dfrac{\pi^2}{L^2}\sin\dfrac{\pi x}{L}-\dfrac{9\pi^2}{4L^2}\sin\dfrac{3\pi x}{L}.
\end{array}
\]

\subsubsection*{Temporal Profile}

For Cases 2a--4a and Problems 7 and 8, the exact solution is
\[
u(x,t)=A\,s(x)g(t),
\qquad
g(t)=1+t^\beta,
\qquad
(x,t)\in(0,L)\times(0,T],
\]
with initial condition $u(x,0)=A\,s(x)$. For $0<\beta<1$, the factor
$\beta t^{\beta-1}$ appearing in the source terms of Cases 2a--4a is
unbounded as $t\to0^+$; the choice $\beta\ge1$ yields a bounded source
term. Problem 5 uses a different temporal profile, given in its
definition.

\subsubsection*{Caputo Derivative}

For $0<\alpha<1$, the Caputo derivative of order $\alpha$ with respect to
$t$ is
\[
\Caputo{\alpha}u(x,t)
=
\frac{1}{\Gamma(1-\alpha)}
\int_0^t (t-\tau)^{-\alpha}
\frac{\partial u(x,\tau)}{\partial \tau}\,d\tau .
\]
It satisfies $\Caputo{\alpha}1=0$ and, for $p>0$,
\[
\Caputo{\alpha}t^p
=
\frac{\Gamma(p+1)}{\Gamma(p+1-\alpha)}\,
t^{p-\alpha}.
\]
For a variable order $\alpha=\alpha(x,t)$ with $0<\alpha(x,t)<1$, we use
the definition in which the order is evaluated at the point $(x,t)$ and
held fixed throughout the integral:
\[
\Caputo{\alpha(x,t)}u(x,t)
=
\frac{1}{\Gamma(1-\alpha(x,t))}
\int_0^t (t-\tau)^{-\alpha(x,t)}
\frac{\partial u(x,\tau)}{\partial \tau}\,d\tau .
\]
Under this definition, the power rule above holds pointwise with
$\alpha$ replaced by $\alpha(x,t)$.

%=====================================================================
\subsection{Elliptic Problems}
%=====================================================================

\subsubsection*{Problem 1. Poisson Equation}

\[
\boxed{-D\,u_{xx}=f(x)}
\]

\paragraph{Case 1a (manufactured).}
$x\in(0,L)$, $u(0)=u(L)=0$.
\[
u(x)=A\,s(x),
\qquad
f(x)=-D\,A\,s''(x).
\]

\paragraph{Case 1b (multiscale solution; Wang, Wang, and Perdikaris \cite{WWP21}).}
$x\in(0,1)$, $u(0)=u(1)=0$.
\[
u(x)=\sin(2\pi x)+0.1\sin(50\pi x),
\qquad
f(x)=D\left[4\pi^2\sin(2\pi x)+250\pi^2\sin(50\pi x)\right].
\]
In \cite{WWP21} the equation is written as $u_{xx}=f$ with $D=1$; the
source term above corresponds to the form $-D\,u_{xx}=f$. The two
components of $f$ have comparable amplitudes ($4\pi^2$ and $250\pi^2$),
whereas the high-frequency component of $u$ is ten times smaller than
the low-frequency one. The case is used to assess the ability of a
method to resolve multiscale solutions.

\paragraph{Case 1c (interior layer).}
$x\in(0,1)$, with Dirichlet data $u(0)$ and $u(1)$ taken from the exact
solution. For $\alpha>0$ and $x_0\in(0,1)$,
\[
u(x)=\arctan\bigl(\alpha(x-x_0)\bigr),
\qquad
f(x)=\frac{2D\alpha^3(x-x_0)}{\bigl[1+\alpha^2(x-x_0)^2\bigr]^2}.
\]
This is the one-dimensional analogue of the ``wave front'' problem in
Mitchell's collection of elliptic test problems \cite{Mit13}, where the
solution is $\arctan(\alpha(r-r_0))$ with $r$ the distance from a fixed
point. The layer at $x=x_0$ has width $O(1/\alpha)$; the values
$\alpha=20$ and $\alpha=1000$ correspond to the mild and steep variants
in \cite{Mit13}. We take $x_0=0.5$.

%=====================================================================
\subsection{Parabolic Problems}
%=====================================================================

\subsubsection*{Problem 2. Heat Equation}

\[
\boxed{u_t-D\,u_{xx}=f(x,t)}
\]

\paragraph{Case 2a (manufactured).}
$(x,t)\in(0,L)\times(0,T]$, $u(0,t)=u(L,t)=0$.
\[
\begin{aligned}
u(x,t)&=A\,s(x)g(t),\\
f(x,t)&=A\left[
\beta t^{\beta-1}s(x)-D\,g(t)s''(x)
\right].
\end{aligned}
\]

\paragraph{Case 2b (free decay of the first eigenmode).}
$(x,t)\in(0,L)\times(0,T]$, $u(0,t)=u(L,t)=0$,
$u(x,0)=\sin(\pi x/L)$, and $f\equiv0$. Separation of variables (see,
e.g., \cite[Ch.~4]{Str08}) gives
\[
u(x,t)=e^{-\lambda t}\sin(\pi x/L),
\qquad
\lambda=\frac{D\pi^2}{L^2}.
\]

\paragraph{Case 2c (time-periodic source).}
$(x,t)\in(0,L)\times(0,T]$, $u(0,t)=u(L,t)=0$, $u(x,0)=0$, and, for a
forcing frequency $\gamma>0$,
\[
f(x,t)=\sin(\pi x/L)\cos(\gamma t).
\]
Since $f$ is proportional to the first eigenfunction, the method of
eigenfunction expansion (see, e.g., \cite[Ch.~8]{Hab13}) reduces the
problem to the scalar equation $a'+\lambda a=\cos(\gamma t)$, $a(0)=0$,
with $\lambda=D\pi^2/L^2$. The solution is
\[
u(x,t)=a(t)\sin(\pi x/L),
\qquad
a(t)=\frac{\lambda\cos(\gamma t)+\gamma\sin(\gamma t)-\lambda e^{-\lambda t}}
{\lambda^2+\gamma^2}.
\]
After the transient $e^{-\lambda t}$ has decayed, the solution
oscillates with amplitude $(\lambda^2+\gamma^2)^{-1/2}$ and a phase lag
of $\arctan(\gamma/\lambda)$ relative to the source.

\subsubsection*{Problem 3. Convection--Diffusion Equation}

\[
\boxed{u_t+v\,u_x-D\,u_{xx}=f(x,t)},
\qquad v\in\mathbb{R}.
\]
The P\'eclet number is $\mathrm{Pe}=vL/D$.

\paragraph{Case 3a (manufactured).}
$(x,t)\in(0,L)\times(0,T]$, $u(0,t)=u(L,t)=0$.
\[
\begin{aligned}
u(x,t)&=A\,s(x)g(t),\\
f(x,t)&=A\left[
\beta t^{\beta-1}s(x)
+v\,g(t)s'(x)-D\,g(t)s''(x)
\right].
\end{aligned}
\]

\paragraph{Case 3b (steady boundary layer, no source; Patankar \cite[Ch.~5]{Pat80}).}
$x\in(0,L)$, $u(0)=0$, $u(L)=1$, $f\equiv0$, $v\neq0$. The
time-independent solution is
\[
u(x)=\frac{e^{\mathrm{Pe}\,x/L}-1}{e^{\mathrm{Pe}}-1}.
\]
For $v>0$ and $\mathrm{Pe}\gg1$, $u$ is close to $0$ except in a boundary layer of
width $O(D/v)$ at $x=L$. This solution is the basis of the exponential
differencing scheme in \cite{Pat80}.

\paragraph{Case 3c (steady boundary layer, constant source).}
$x\in(0,L)$, $u(0)=u(L)=0$, $f\equiv1$, $v\neq0$. Direct integration of
$v\,u'-D\,u''=1$ gives
\[
u(x)=\frac{1}{v}\left[x-L\,\frac{e^{\mathrm{Pe}\,x/L}-1}{e^{\mathrm{Pe}}-1}\right].
\]
For $v>0$ and $\mathrm{Pe}\gg1$, $u\approx x/v$ away from $x=L$, with a
boundary layer of width $O(D/v)$ at $x=L$. This is a standard model
problem for singularly perturbed convection--diffusion equations; see
\cite{RST08}.

\paragraph{Case 3d (advected Gaussian pulse).}
$x\in\mathbb{R}$, $f\equiv0$. For $\sigma_0>0$ and $x_0\in\mathbb{R}$,
\[
u(x,t)=\frac{\sigma_0}{\sigma(t)}
\exp\!\left(-\frac{(x-x_0-vt)^2}{2\sigma(t)^2}\right),
\qquad
\sigma(t)^2=\sigma_0^2+2Dt.
\]
This follows from the fundamental solution of the heat equation
\cite[Sec.~2.3]{Eva10} and the change of variables $\xi=x-vt$. On a
bounded interval, the Dirichlet data are taken from the exact solution;
if the interval is chosen so that the pulse remains far from the
boundary, these data are negligibly small.

\subsubsection*{Problem 4. Viscous Burgers Equation}

The equation is quasilinear and, since $D>0$, parabolic.
\[
\boxed{u_t+u\,u_x-D\,u_{xx}=f(x,t)}
\]

\paragraph{Case 4a (manufactured).}
$(x,t)\in(0,L)\times(0,T]$, $u(0,t)=u(L,t)=0$.
\[
\begin{aligned}
u(x,t)&=A\,s(x)g(t),\\
f(x,t)&=A\left[
\beta t^{\beta-1}s(x)-D\,g(t)s''(x)
\right]
+A^2g(t)^2s(x)s'(x).
\end{aligned}
\]

\paragraph{Case 4b (sine initial data; Basdevant et al.\ \cite{Bas86}).}
$(x,t)\in(-1,1)\times(0,1]$, $u(\pm1,t)=0$, $u(x,0)=-\sin(\pi x)$,
$f\equiv0$, and $D=0.01/\pi$. The Cole--Hopf transformation gives
\[
u(x,t)=-\,\frac{\displaystyle\int_{-\infty}^{\infty}
\sin\bigl(\pi(x-\eta)\bigr)\,h(x-\eta)\,e^{-\eta^2/(4Dt)}\,d\eta}
{\displaystyle\int_{-\infty}^{\infty}
h(x-\eta)\,e^{-\eta^2/(4Dt)}\,d\eta},
\qquad
h(y)=\exp\!\left(-\frac{\cos(\pi y)}{2\pi D}\right).
\]
With the substitution $\eta=2\sqrt{Dt}\,z$, both integrals take the
form $\int_{-\infty}^{\infty}e^{-z^2}(\cdot)\,dz$ and can be evaluated
by Gauss--Hermite quadrature, as in \cite{Bas86}. The solution develops
a steep internal layer at $x=0$. This case was used as a benchmark for
physics-informed neural networks by Raissi, Perdikaris, and Karniadakis
\cite{RPK19}.

\paragraph{Case 4c (viscous shock; Whitham \cite[Ch.~4]{Whi74}).}
$x\in\mathbb{R}$, $f\equiv0$. For constants $u_\ell>u_r$, the travelling
wave
\[
u(x,t)=\frac{u_\ell+u_r}{2}
-\frac{u_\ell-u_r}{2}
\tanh\!\left(\frac{(u_\ell-u_r)(x-ct)}{4D}\right),
\qquad
c=\frac{u_\ell+u_r}{2},
\]
connects $u_\ell$ as $x\to-\infty$ to $u_r$ as $x\to+\infty$ and moves
with the Rankine--Hugoniot speed $c$. The shock width is
$O\bigl(D/(u_\ell-u_r)\bigr)$. On a bounded interval, the Dirichlet data
are taken from the exact solution.

%=====================================================================
\subsection{Hyperbolic Problems}
%=====================================================================

\subsubsection*{Problem 5. Two-Dimensional Wave Equation}

\[
\boxed{u_{tt}-D(u_{xx}+u_{yy})=f(x,y,t)},
\qquad (x,y)\in\Omega=(0,L)^2,
\quad t\in(0,T],
\quad u=0 \text{ on } \partial\Omega .
\]
Define
\[
S(x,y)=s(x)s(y),
\qquad
\Delta S=s''(x)s(y)+s(x)s''(y),
\qquad
\omega=\frac{\pi\sqrt{2D}}{L}.
\]
Exact solution and initial conditions:
\[
\begin{aligned}
u(x,y,t)&=A\,S(x,y)\left[\cos(\omega t)+t^\beta\right],\\
u(x,y,0)&=A\,S(x,y),\\
u_t(x,y,0)&=0.
\end{aligned}
\]
Source term:
\[
f(x,y,t)
=A\left[
-\omega^2\cos(\omega t)
+\beta(\beta-1)t^{\beta-2}
\right]S(x,y)
-DA\left[\cos(\omega t)+t^\beta\right]\Delta S(x,y).
\]
The initial condition $u_t(x,y,0)=0$ requires $\beta>1$. For
$1<\beta<2$, the term $\beta(\beta-1)t^{\beta-2}$ is unbounded as
$t\to0^+$; we therefore take $\beta\ge2$.

For the training profile, $S$ is the first Dirichlet eigenfunction of
$-\Delta$ on $\Omega$, with eigenvalue $\lambda_1=2\pi^2/L^2$, and
$\omega=\sqrt{D\lambda_1}$. Consequently, $A\,S\cos(\omega t)$ solves the
homogeneous equation, and the source term reduces to
\[
f(x,y,t)=A\left[\beta(\beta-1)t^{\beta-2}+\frac{2\pi^2D}{L^2}\,t^\beta\right]S(x,y).
\]
For the validation and reserved profiles, $S$ is not an eigenfunction,
and no such reduction occurs.

%=====================================================================
\subsection{Time-Fractional Problems}
%=====================================================================

The classification into elliptic, parabolic, and hyperbolic types is
defined for local differential operators through their principal part.
The Caputo derivative is a nonlocal operator in time, so this
classification does not apply to the problems below. For $0<\alpha<1$,
the problems in this section are time-fractional diffusion equations,
also called subdiffusion equations; they reduce formally to the heat
equation as $\alpha\to1^-$. For $1<\alpha<2$, the corresponding
equation is the diffusion-wave equation, which interpolates between the
heat equation ($\alpha=1$) and the wave equation ($\alpha=2$); no such
problem is included here.

\subsubsection*{Scalar Problem}

\subsubsection*{Problem 6. Fractional Ordinary Differential Equation}

\[
\boxed{\Caputo{\alpha}u(t)=f(t)},
\qquad t\in(0,1],
\qquad u(0)=0,
\qquad 0<\alpha<1,
\quad p>0.
\]
Exact solution and source term:
\[
u(t)=t^p,
\qquad
f(t)=
\frac{\Gamma(p+1)}{\Gamma(p+1-\alpha)}\,
t^{p-\alpha}.
\]
The source term is bounded on $(0,1]$ if and only if $p\ge\alpha$.

\subsubsection*{Constant-Order Time-Fractional Diffusion}

\subsubsection*{Problem 7. Time-Fractional Diffusion Equation}

\[
\boxed{\Caputo{\alpha}u-D\,u_{xx}=f(x,t)},
\qquad 0<\alpha<1.
\]
Exact solution and source term:
\[
\begin{aligned}
u(x,t)&=A\,s(x)g(t),\\
f(x,t)&=A\left[
\frac{\Gamma(\beta+1)}{\Gamma(\beta+1-\alpha)}\,
t^{\beta-\alpha}s(x)-D\,g(t)s''(x)
\right].
\end{aligned}
\]
The source term is bounded on $(0,L)\times(0,T]$ if and only if
$\beta\ge\alpha$.

\subsubsection*{Problem 8. Multi-Term Time-Fractional Diffusion Equation}

\[
\boxed{
\sum_{j=1}^{2}w_j\,\Caputo{\alpha_j}u
-D\,u_{xx}=f(x,t)
},
\qquad 0<\alpha_j<1,
\quad w_j>0.
\]
Exact solution and source term:
\[
\begin{aligned}
u(x,t)&=A\,s(x)g(t),\\
f(x,t)&=
A\left[
\sum_{j=1}^{2}
w_j\frac{\Gamma(\beta+1)}
{\Gamma(\beta+1-\alpha_j)}\,
t^{\beta-\alpha_j}s(x)
-D\,g(t)s''(x)
\right].
\end{aligned}
\]
The source term is bounded on $(0,L)\times(0,T]$ if and only if
$\beta\ge\max(\alpha_1,\alpha_2)$.

\subsubsection*{Problem 9. Time-Fractional Diffusion Equation with Zero Initial Data}

Problems 9--11 are posed on $(x,t)\in(0,\pi)\times(0,1]$ with $p>0$,
boundary conditions $u(0,t)=u(\pi,t)=0$, and initial condition
$u(x,0)=0$.
\[
\boxed{\Caputo{\alpha}u=u_{xx}+f(x,t)},
\qquad 0<\alpha<1.
\]
Exact solution and source term:
\[
\begin{aligned}
u(x,t)&=t^p\sin x,\\
f(x,t)&=
\left[
\frac{\Gamma(p+1)}{\Gamma(p+1-\alpha)}\,
t^{p-\alpha}+t^p
\right]\sin x.
\end{aligned}
\]
This problem coincides with Problem 7 for $D=1$, $L=\pi$, $A=1$, and the
training profile, except that the temporal profile is $t^p$ rather than
$1+t^\beta$.

\subsubsection*{Variable-Order Time-Fractional Diffusion}

\subsubsection*{Problem 10. Time-Dependent Order}

\[
\boxed{\Caputo{\alpha(t)}u=u_{xx}+f(x,t)},
\qquad \alpha(t)=c\sin^2 t,
\qquad 0<c<\frac{1}{\sin^2 1}\approx1.412.
\]
Exact solution and source term:
\[
\begin{aligned}
u(x,t)&=t^p\sin x,\\
f(x,t)&=
\left[
\frac{\Gamma(p+1)}{\Gamma(p+1-\alpha(t))}\,
t^{p-\alpha(t)}+t^p
\right]\sin x.
\end{aligned}
\]
Since $\sin^2 t$ is positive and increasing on $(0,1]$, the stated bound
on $c$ is equivalent to $0<\alpha(t)<1$ for all $t\in(0,1]$. Note that
$\alpha(t)\to0$ as $t\to0^+$.

\subsubsection*{Problem 11. Space- and Time-Dependent Order}

\[
\boxed{\Caputo{\alpha(x,t)}u=u_{xx}+f(x,t)},
\qquad \alpha(x,t)=c\,t^2\sin x,
\qquad 0<c<1.
\]
Exact solution and source term:
\[
\begin{aligned}
u(x,t)&=t^p\sin x,\\
f(x,t)&=
\left[
\frac{\Gamma(p+1)}{\Gamma(p+1-\alpha(x,t))}\,
t^{p-\alpha(x,t)}+t^p
\right]\sin x.
\end{aligned}
\]
Since $\max_{(0,\pi)\times(0,1]}t^2\sin x=1$, attained at
$(x,t)=(\pi/2,1)$, the stated bound on $c$ is equivalent to
$0<\alpha(x,t)<1$ on $(0,\pi)\times(0,1]$. Note that $\alpha(x,t)\to0$
as $t\to0^+$ or as $x$ approaches $0$ or $\pi$.

%=====================================================================
\subsection{Fourth-Order and Incompressible Flow Problems}
%=====================================================================

\subsubsection*{Problem 12. Kuramoto--Sivashinsky Equation}

\[
\boxed{u_t+u\,u_x+u_{xx}+u_{xxxx}=f(x,t)},
\qquad x\in(0,L),
\quad t\in(0,T].
\]
Boundary conditions are periodic:
\[
\frac{\partial^j u}{\partial x^j}(0,t)=\frac{\partial^j u}{\partial x^j}(L,t),
\qquad j=0,1,2,3,
\quad t\in(0,T].
\]
The term $u_{xx}$ is destabilizing and $u_{xxxx}$ is stabilizing. The
Fourier mode $e^{ikx}$ grows under the linear part if and only if
$0<|k|<1$.

\paragraph{Case 12a (manufactured).}
Let $k=2\pi/L$. Then
\[
\begin{aligned}
u(x,t)&=A\sin(kx)\,g(t),\\
f(x,t)&=A\left[\beta t^{\beta-1}+(k^4-k^2)\,g(t)\right]\sin(kx)
+\frac{A^2k}{2}\,g(t)^2\sin(2kx),
\end{aligned}
\]
with initial condition $u(x,0)=A\sin(kx)$. As in Cases 2a--4a, the
source term is bounded if and only if $\beta\ge1$.

\paragraph{Case 12b (chaotic dynamics; Kassam and Trefethen \cite{KT05}).}
$L=32\pi$, $T=150$, $f\equiv0$, and
\[
u(x,0)=\cos\!\left(\frac{x}{16}\right)\left[1+\sin\!\left(\frac{x}{16}\right)\right].
\]
There is no closed-form solution. A reference solution is computed by a
Fourier spectral method in space with the ETDRK4 scheme in time, as in
\cite{KT05}. The wavenumbers are $n/16$, $n\in\mathbb{Z}$, so the
modes $1\le|n|\le15$ are linearly unstable. The solution evolves into
spatiotemporal chaos.

\subsubsection*{Problem 13. Two-Dimensional Incompressible Navier--Stokes Equations}

\[
\boxed{
\mathbf{u}_t+(\mathbf{u}\cdot\nabla)\mathbf{u}+\nabla p-\nu\,\Delta\mathbf{u}=\mathbf{f},
\qquad
\nabla\cdot\mathbf{u}=0
},
\]
where $\mathbf{u}=(u,v)$ is the velocity, $p$ is the pressure (unit
density), and $\nu>0$ is the kinematic viscosity. The Reynolds number is
$\mathrm{Re}=1/\nu$.

\paragraph{Case 13a (Taylor--Green vortex; Taylor and Green \cite{TG37}).}
$(x,y)\in\Omega=(0,2\pi)^2$, $t\in(0,T]$, $\mathbf{f}\equiv\mathbf{0}$.
Boundary conditions are periodic in both directions for $\mathbf{u}$ and $p$:
\[
\mathbf{u}(0,y,t)=\mathbf{u}(2\pi,y,t),\qquad
\mathbf{u}(x,0,t)=\mathbf{u}(x,2\pi,t),
\]
and likewise for $p$ and the first derivatives of $\mathbf{u}$. The
initial condition is
\[
\mathbf{u}(x,y,0)=(\cos x\,\sin y,\;-\sin x\,\cos y).
\]
The exact solution is
\[
\begin{aligned}
u(x,y,t)&=\cos x\,\sin y\;F(t),\\
v(x,y,t)&=-\sin x\,\cos y\;F(t),\\
p(x,y,t)&=-\frac14\bigl[\cos(2x)+\cos(2y)\bigr]F(t)^2,
\end{aligned}
\qquad F(t)=e^{-2\nu t}.
\]
The pressure is determined up to an additive constant, which is fixed by
$\int_\Omega p\,dx\,dy=0$; the exact pressure satisfies this condition.
The pressure gradient balances the convective term exactly, and
$\Delta\mathbf{u}=-2\mathbf{u}$. The flow therefore decays without
changing shape. The benchmark also uses the vortex of integer wavenumber
$k$, with $\cos x,\sin y,\ldots$ replaced by $\cos kx,\sin ky,\ldots$,
$\cos 2x,\cos 2y$ by $\cos 2kx,\cos 2ky$, and $F(t)=e^{-2\nu k^2t}$.

\paragraph{Case 13b (Kovasznay flow; Kovasznay \cite{Kov48}).}
$(x,y)\in\Omega=(-0.5,1)\times(-0.5,1.5)$, steady,
$\mathbf{f}\equiv\mathbf{0}$. Let
\[
\lambda=\frac{\mathrm{Re}}{2}-\sqrt{\frac{\mathrm{Re}^2}{4}+4\pi^2}.
\]
The exact solution is
\[
\begin{aligned}
u(x,y)&=1-e^{\lambda x}\cos(2\pi y),\\
v(x,y)&=\frac{\lambda}{2\pi}\,e^{\lambda x}\sin(2\pi y),\\
p(x,y)&=\frac12\bigl(1-e^{2\lambda x}\bigr).
\end{aligned}
\]
The velocity satisfies the Dirichlet condition $\mathbf{u}=(u,v)$ on
$\partial\Omega$, with $(u,v)$ given above. Since this velocity is
divergence-free, the compatibility condition
$\int_{\partial\Omega}\mathbf{u}\cdot\mathbf{n}\,ds=0$ holds. No boundary
condition is imposed on $p$. The pressure is determined up to an
additive constant, which is fixed by matching the exact pressure at one
point, e.g.\ $p(-0.5,-0.5)=\tfrac12\bigl(1-e^{-\lambda}\bigr)$. This
solution models laminar flow behind a two-dimensional grid. A typical
choice is $\mathrm{Re}=40$.

%=====================================================================
\subsection{Phase-Field and Electromagnetic Problems}
%=====================================================================

\subsubsection*{Problem 14. Cahn--Hilliard Equation}

\[
\boxed{u_t=M\,\Delta\mu+f(x,y,t),\qquad \mu=u^3-u-\varepsilon^2\Delta u},
\qquad (x,y)\in(0,L)^2,
\quad t\in(0,T],
\]
with mobility $M>0$, interface parameter $\varepsilon>0$, and periodic
boundary conditions in both directions \cite{CH58}. Without forcing the
equation is the $H^{-1}$ gradient flow of the free energy
$\int\bigl[\tfrac14(u^2-1)^2+\tfrac{\varepsilon^2}{2}|\nabla u|^2\bigr]$,
so $\int_\Omega u$ is conserved. Linearized about $u=0$, the Fourier mode
of wavenumber $|\mathbf{k}|$ grows at the rate
$M|\mathbf{k}|^2(1-\varepsilon^2|\mathbf{k}|^2)$, so the modes with
$0<|\mathbf{k}|<1/\varepsilon$ are unstable.

\paragraph{Case 14a (manufactured).}
Let $k=2\pi/L$, $c(x,y)=\cos(kx)\cos(ky)$, and $g(t)=1+t^\beta$. Then
\[
\begin{aligned}
u(x,y,t)&=A\,c(x,y)\,g(t),\\
f(x,y,t)&=A\,\beta t^{\beta-1}c(x,y)
-M\left[3u^2\Delta u+6u\,|\nabla u|^2-\Delta u-4\varepsilon^2k^4u\right],
\end{aligned}
\]
where $\Delta u=-2k^2u$ and
$|\nabla u|^2=A^2g(t)^2k^2\bigl[\sin^2(kx)\cos^2(ky)+\cos^2(kx)\sin^2(ky)\bigr]$.
The identity $\Delta(u^3)=3u^2\Delta u+6u|\nabla u|^2$ gives the
nonlinear part. The source term is bounded if and only if $\beta\ge1$.

\paragraph{Case 14b (spinodal decomposition).}
$L=2\pi$, $M=1$, $f\equiv0$, and a small smooth perturbation of the
mixed state,
\[
u(x,y,0)=\sum_{j=1}^{5}a_j\cos\bigl(k^{(j)}_xx+k^{(j)}_yy+\phi_j\bigr),
\]
with amplitudes $a_j\in[0.04,0.1]$, integer wavenumbers, and phases that
differ between the splits, as does $\varepsilon\in\{0.15,0.12,0.1\}$.
The unstable modes grow until the field separates into regions with
$u\approx\pm1$ divided by interfaces of width $O(\varepsilon)$. There is
no closed-form solution. A reference solution is computed by a Fourier
spectral method with the ETDRK4 scheme \cite{KT05}, on $256^2$ modes, and
is checked against a $128^2$ solve.

\subsubsection*{Problem 15. Maxwell's Equations (Two-Dimensional TE Mode)}

In units in which the permittivity and permeability are equal and
$c=1/\sqrt{\epsilon\mu}$ is the wave speed, the time-domain Maxwell
equations for the transverse electric field
$\mathbf{E}=(E_x,E_y)$ and the out-of-plane magnetic field $H_z$ read
\cite{Jac99,Yee66}
\[
\boxed{
\mathbf{E}_t=c\,\nabla\times H_z-\mathbf{J},
\qquad
(H_z)_t=-c\,\nabla\times\mathbf{E}
},
\]
with $\nabla\times H_z=(\partial_yH_z,\,-\partial_xH_z)$ and
$\nabla\times\mathbf{E}=\partial_xE_y-\partial_yE_x$, on
$\Omega=(0,a)\times(0,b)$ and $t\in(0,T]$. The boundary is a perfect
electric conductor: the tangential electric field vanishes,
$E_x=0$ on $y=0,b$ and $E_y=0$ on $x=0,a$.

\paragraph{Case 15a (resonant cavity mode).}
$\mathbf{J}\equiv\mathbf{0}$. For integers $m,n\ge1$, let
$\omega=c\pi\sqrt{(m/a)^2+(n/b)^2}$. The standing wave
\[
\begin{aligned}
H_z&=\cos\frac{m\pi x}{a}\cos\frac{n\pi y}{b}\cos\omega t,\\
E_x&=-\frac{cn\pi}{b\omega}\cos\frac{m\pi x}{a}\sin\frac{n\pi y}{b}\sin\omega t,\\
E_y&=\frac{cm\pi}{a\omega}\sin\frac{m\pi x}{a}\cos\frac{n\pi y}{b}\sin\omega t
\end{aligned}
\]
is the TE$_{mn}$ eigenmode of the cavity, with $\mathbf{E}(x,y,0)=\mathbf{0}$.

\paragraph{Case 15b (current-driven field; manufactured).}
For $A\neq0$ and $\beta\ge1$,
\[
\begin{aligned}
\mathbf{E}&=A\,t^\beta\Bigl(\cos\frac{\pi x}{a}\sin\frac{\pi y}{b},\;
-\sin\frac{\pi x}{a}\cos\frac{\pi y}{b}\Bigr),\\
H_z&=cA\pi\Bigl(\frac1a+\frac1b\Bigr)\frac{t^{\beta+1}}{\beta+1}
\cos\frac{\pi x}{a}\cos\frac{\pi y}{b},
\end{aligned}
\]
which satisfies Faraday's law exactly and vanishes at $t=0$. The current
density is $\mathbf{J}=c\,\nabla\times H_z-\mathbf{E}_t$.

\subsubsection*{Problem 16. Time-Harmonic Electromagnetic Equation}

For a field oscillating as $e^{-i\omega t}$ in a lossless medium, Maxwell's
equations reduce to the curl--curl equation \cite{Mon03}
\[
\boxed{\nabla\times\nabla\times\mathbf{E}-k^2\mathbf{E}=\mathbf{f}},
\qquad (x,y)\in\Omega=(0,1)^2,
\qquad \mathbf{n}\times\mathbf{E}=0 \text{ on } \partial\Omega,
\]
with wavenumber $k>0$ and, in two dimensions,
$\nabla\times\nabla\times\mathbf{E}=(\partial_yw,\,-\partial_xw)$,
$w=\partial_xE_y-\partial_yE_x$. The problem is uniquely solvable when
$k^2$ is not a Maxwell eigenvalue of the square, $\pi^2(m^2+n^2)$ with
integers $m,n\ge0$ not both zero; every case takes $k$ strictly between two
consecutive eigenvalues.

\paragraph{Cases 16a and 16b (manufactured).}
For integers $p,n,m,q\ge1$ and amplitudes $A,B$,
\[
\begin{aligned}
E_x&=A\cos(p\pi x)\sin(n\pi y),\qquad
E_y=B\sin(m\pi x)\cos(q\pi y),\\
f_x&=\pi^2\bigl[n^2A\cos(p\pi x)\sin(n\pi y)-mqB\cos(m\pi x)\sin(q\pi y)\bigr]-k^2E_x,\\
f_y&=\pi^2\bigl[m^2B\sin(m\pi x)\cos(q\pi y)-npA\sin(p\pi x)\cos(n\pi y)\bigr]-k^2E_y.
\end{aligned}
\]
The tangential components vanish on the boundary. Case 16a uses a low
frequency ($k<\pi$). Case 16b uses a high frequency near the resonances,
for example $k=10.6$, between $\pi\sqrt{10}$ and $\pi\sqrt{13}$, together
with higher modes.


\subsection*{Numbers of the public cases}
Training (evolution) and validation (development) cases. The literature cases use the values below. The manufactured cases (1a, 2a, 3a, 4a, 5, 7, 8) draw $L,T,A,D,\beta$ (and $v$, $\alpha$, $w_j$) per case and split from the ranges in \texttt{pdebench.manufactured}. The sealed test splits (test: 17 cases, test2 and test3: 20 cases each) use further values that are withheld here.
\begin{center}
\footnotesize
\begin{tabular}{lp{3.6cm}p{4.6cm}p{4.6cm}}
\toprule
Case & Title & Training & Validation \\
\midrule
1b & Poisson, multiscale & --- & $D$ $=$ 1 \\
1c & Poisson, interior layer & $D$ $=$ 1, $\alpha$ $=$ 20, $x_0$ $=$ 0.5 & --- \\
2b & Heat, free decay & $L$ $=$ 2, $D$ $=$ 0.3 & --- \\
2c & Heat, time-periodic source & --- & $T$ $=$ 2, $D$ $=$ 0.2, $\gamma$ $=$ $2\pi$ \\
3b & Convection--diffusion, steady boundary layer & $v$ $=$ 1, $D$ $=$ 0.05 & --- \\
3c & Convection--diffusion, steady layer with constant source & --- & $v$ $=$ 1, $D$ $=$ 0.02 \\
3d & Convection--diffusion, advected Gaussian & $L$ $=$ 4, $v$ $=$ 1, $D$ $=$ 0.02, $x_0$ $=$ 1, $\sigma_0$ $=$ 0.2 & $L$ $=$ 4, $v$ $=$ -0.8, $D$ $=$ 0.015, $x_0$ $=$ 2.5, $\sigma_0$ $=$ 0.15 \\
4b & Burgers, sine data & --- & $D=0.01/\pi$ \\
4c & Burgers, viscous shock & $L$ $=$ 1, $u_\ell$ $=$ 1, $u_r$ $=$ 0, $D$ $=$ 0.02 & --- \\
6 & Fractional ODE & $\alpha$ $=$ 0.4, $p$ $=$ 2 & $\alpha$ $=$ 0.7, $p$ $=$ 0.8 \\
9 & Time-fractional diffusion, zero initial data & $\alpha$ $=$ 0.5, $p$ $=$ 2 & $\alpha$ $=$ 0.8, $p$ $=$ 1.5 \\
10 & Variable order alpha(t) & $c$ $=$ 0.8, $p$ $=$ 2 & $c$ $=$ 1.3, $p$ $=$ 2.5 \\
11 & Variable order alpha(x,t) & $c$ $=$ 0.5, $p$ $=$ 2 & $c$ $=$ 0.9, $p$ $=$ 3 \\
\bottomrule
\end{tabular}
\end{center}
\subsection*{Numbers of the public cases, Problems 12--16}
\begin{center}
\footnotesize
\begin{tabular}{lp{3.4cm}p{4.7cm}p{4.7cm}}
\toprule
Case & Title & Training & Validation \\
\midrule
12a & Kuramoto--Sivashinsky, manufactured & $L$ $=$ 10.52, $A$ $=$ 1.465, $\beta$ $=$ 1.175, $T$ $=$ 0.5839 & $L$ $=$ 15.19, $A$ $=$ 0.8677, $\beta$ $=$ 1.257, $T$ $=$ 0.9132 \\
12b & Kuramoto--Sivashinsky, chaotic (Kassam--Trefethen) & $L$ $=$ $32\pi$, $T$ $=$ 30 & $L$ $=$ $32\pi$, $T$ $=$ 60 \\
13a & Navier--Stokes, Taylor--Green vortex & $k$ $=$ 1, $\nu$ $=$ 0.1, $T$ $=$ 1 & $k$ $=$ 2, $\nu$ $=$ 0.05, $T$ $=$ 1 \\
13b & Navier--Stokes, Kovasznay flow & $\nu$ $=$ 0.05 & $\nu$ $=$ 0.025 \\
14a & Cahn--Hilliard, manufactured & $L$ $=$ 3.263, $A$ $=$ 0.5148, $\beta$ $=$ 1.762, $M$ $=$ 1.382, $\varepsilon$ $=$ 0.06735, $T$ $=$ 0.2427 & $L$ $=$ 2.771, $A$ $=$ 0.306, $\beta$ $=$ 1.976, $M$ $=$ 1.326, $\varepsilon$ $=$ 0.09144, $T$ $=$ 0.2588 \\
14b & Cahn--Hilliard, spinodal decomposition & $\varepsilon$ $=$ 0.15, $M$ $=$ 1, $T$ $=$ 0.5 & $\varepsilon$ $=$ 0.12, $M$ $=$ 1, $T$ $=$ 0.4 \\
15a & Maxwell TE$_z$, resonant cavity mode & $a$ $=$ 1, $b$ $=$ 1, $m$ $=$ 1, $n$ $=$ 1, $c$ $=$ 1, $T$ $=$ 1 & $a$ $=$ 1, $b$ $=$ 1.5, $m$ $=$ 2, $n$ $=$ 1, $c$ $=$ 1.5, $T$ $=$ 1 \\
15b & Maxwell TE$_z$, current-driven (manufactured) & $a$ $=$ 1.992, $b$ $=$ 1.984, $A$ $=$ 0.8852, $\beta$ $=$ 1.292, $c$ $=$ 1.881, $T$ $=$ 0.6036 & $a$ $=$ 1.503, $b$ $=$ 1.301, $A$ $=$ 0.878, $\beta$ $=$ 1.036, $c$ $=$ 1.897, $T$ $=$ 0.9697 \\
16a & Time-harmonic electromagnetic, low frequency & $A$ $=$ 1, $B$ $=$ 0.5, $p$ $=$ 1, $n$ $=$ 1, $m$ $=$ 1, $q$ $=$ 1, $k$ $=$ 2 & $A$ $=$ 0.8, $B$ $=$ 1.2, $p$ $=$ 2, $n$ $=$ 1, $m$ $=$ 1, $q$ $=$ 2, $k$ $=$ 1.5 \\
16b & Time-harmonic electromagnetic, high frequency & $A$ $=$ 1, $B$ $=$ 0.7, $p$ $=$ 2, $n$ $=$ 3, $m$ $=$ 3, $q$ $=$ 1, $k$ $=$ 10.6 & $A$ $=$ 0.9, $B$ $=$ 1.1, $p$ $=$ 1, $n$ $=$ 2, $m$ $=$ 2, $q$ $=$ 3, $k$ $=$ 7.6 \\
\bottomrule
\end{tabular}
\end{center}
The sealed test split of Problems 12--16 has one case per row above (10 cases); its numbers are withheld.
\section{PDE Data Inventory suite}

Twelve families from \emph{PDE Data Inventory} (X.~Kong, 19 October 2026), with inputs from the DiffusionPDE test
set (Huang et al.\ 2024) where the inventory leaves gaps. The library adds the following classes to the ten classes
of Problems 1--11:
\begin{align*}
&\text{darcy\_2d}:\quad -\nabla\cdot\bigl(a(x,y)\nabla u\bigr)=f(x,y)\\
&\text{poisson\_2d}:\quad u_{xx}+u_{yy}=f(x,y)\\
&\text{helmholtz\_2d}:\quad u_{xx}+u_{yy}+k^2u=f(x,y)\\
&\text{vorticity\_ns}:\quad \omega_t+\mathbf{v}\cdot\nabla\omega=\nu\Delta\omega-\mu\omega+q,\quad \mathbf{v}=(\psi_y,-\psi_x),\ -\Delta\psi=\omega\\
&\text{allen\_cahn\_2d}:\quad u_t=\nabla\cdot(\kappa\nabla u)+\lambda(u-u^3)+f\\
&\text{burgers\_2d}:\quad u_t+c_xuu_x+c_yuu_y=\nabla\cdot(\kappa\nabla u)+f\\
&\text{adr\_2d}:\quad u_t+c_xuu_x+c_yuu_y=\nabla\cdot(\kappa\nabla u)+\lambda(u-u^3)+f\\
&\text{elliptic\_4d}:\quad -\nabla\cdot(a\nabla u)=f \text{ in } B^3\times(-1,1),\quad u=0 \text{ on the boundary}
\end{align*}
\begin{center}
\footnotesize\setlength{\tabcolsep}{3pt}
\begin{tabular}{p{2.4cm}p{2.9cm}p{2.45cm}p{0.8cm}p{0.9cm}p{2.4cm}p{2.4cm}}
\toprule
Family & Class & Layout & Limit & Tol. & Training & Validation \\
\midrule
1-D Burgers, fixed Gaussian initial state & burgers & 1d\_periodic & 90 s & 0.002 & $D=0.1$ & $D=0.07$ \\
1-D Burgers, random periodic initial state & burgers & 1d\_periodic & 90 s & 0.002 & instance 0 & instance 1 \\
1-D linear convection--diffusion & convection\_diffusion & 1d\_dirichlet & 90 s & 0.002 & $x_0=5$, $\sigma_0=1$ & $x_0=4$, $\sigma_0=0.8$ \\
Darcy flow, binary permeability & darcy\_2d & 2d\_steady & 300 s & 0.002 & instance 0 & instance 1 \\
Poisson equation, random forcing & poisson\_2d & 2d\_steady & 300 s & 0.002 & instance 0 & instance 1 \\
Inhomogeneous Helmholtz, k = 1 & helmholtz\_2d & 2d\_steady & 300 s & 0.002 & instance 0 & instance 1 \\
2-D Navier--Stokes vorticity, unit torus & vorticity\_ns & 2d\_periodic\_time & 300 s & 0.002 & instance 0 & instance 1 \\
2-D Kolmogorov flow with drag & vorticity\_ns & 2d\_periodic\_time & 300 s & 0.002 & own draw & own draw \\
Heterogeneous Allen--Cahn & allen\_cahn\_2d & 2d\_periodic\_time & 300 s & 0.002 & $\lambda=0.840$ & $\lambda=0.833$ \\
2-D stochastic scalar Burgers & burgers\_2d & 2d\_periodic\_time & 300 s & 0.002 & $c_x=0.813$, $c_y=0.546$ & $c_x=1.020$, $c_y=0.504$ \\
Stochastic nonlinear advection--diffusion--reaction & adr\_2d & 2d\_periodic\_time & 300 s & 0.002 & $c_x=0.617$, $c_y=0.255$, $\lambda=0.520$ & $c_x=0.644$, $c_y=0.378$, $\lambda=0.375$ \\
4-D elliptic, invariant coefficient & elliptic\_4d & 4d\_points & 300 s & 0.05 & own draw & own draw \\
\bottomrule
\end{tabular}
\end{center}
\noindent The sealed splits \texttt{inventory\_test} and \texttt{inventory\_test2} (DiffusionPDE
instances 2 and 3, and further draws) are withheld. References are computed at high resolution from the same
continuous inputs: pseudospectral ETDRK4 (1-D Burgers), finite volumes on an $8\times$ refined grid (Darcy), sine
series (Poisson, Helmholtz), integrating-factor RK4 (Navier--Stokes, Kolmogorov), pseudospectral RK4 (Allen--Cahn,
2-D Burgers, ADR), $63^4$ finite volumes (4-D elliptic) and the closed form (linear convection--diffusion). Each is
checked against a coarser solve; the tolerance is at least five times the reference error.
\section{Manufactured suites (code protocol)}\label{sec:manufactured}
\renewcommand{\thesubsection}{\thesection.\arabic{subsection}}
Generated from \texttt{src/pde\_rsi/tasks.py} by \texttt{scripts/write\_problem\_latex.py}. Every problem has a manufactured exact solution; agents see the forcing and initial data only as black boxes.

\subsection{One-dimensional benchmark (48 public, 24 reserved cases)}
All cases use $x\in[0,L]$, $u(0,t)=u(L,t)=0$ and the exact solution
\begin{equation*}
u(x,t) = A\,s(x)\,\bigl(1+t^{\beta}\bigr) \quad (\text{Poisson: } u(x)=A\,s(x)), \qquad u(x,0)=A\,s(x),
\end{equation*}
with the spatial profile fixed by the split:
\begin{equation*}
s(x)=\begin{cases}
\sin(\pi x/L) & \text{train,}\\
4\,\dfrac{x}{L}\Bigl(1-\dfrac{x}{L}\Bigr) & \text{validation,}\\
\sin(\pi x/L)+\tfrac14\sin(3\pi x/L) & \text{reserved.}
\end{cases}
\end{equation*}
Write $g(t)=1+t^{\beta}$.

\subsubsection{Poisson}
\begin{align*}
-D\,u_{xx} = f(x),\quad x\in(0,L)\\
f = -D\,A\,s''(x)
\end{align*}
\begin{center}\small
\begin{tabular}{llrrrr}
\toprule
Case & Split & $L$ & $D$ & $A$ & Profile \\
\midrule
\texttt{poisson-g0-v0} & train & 1 & 0.05 & 0.3 & sine \\
\texttt{poisson-g0-v1} & train & 1.2 & 0.1 & 0.4 & sine \\
\texttt{poisson-g0-v2} & train & 1 & 0.15 & 0.5 & sine \\
\texttt{poisson-g0-v3} & train & 1.2 & 0.2 & 0.6 & sine \\
\texttt{poisson-g1-v0} & validation & 1 & 0.05 & 0.3 & quadratic \\
\texttt{poisson-g1-v1} & validation & 1.2 & 0.1 & 0.4 & quadratic \\
\texttt{poisson-g1-v2} & validation & 1 & 0.15 & 0.5 & quadratic \\
\texttt{poisson-g1-v3} & validation & 1.2 & 0.2 & 0.6 & quadratic \\
\bottomrule
\end{tabular}
\end{center}
\noindent Reserved cases: 4 (parameters kept evaluator-only).

\subsection{Heat}
\begin{align*}
u_t - D\,u_{xx} = f(x,t),\quad x\in(0,L),\ t\in(0,T]\\
f = A\bigl[\beta t^{\beta-1}s(x) - D\,g(t)\,s''(x)\bigr]
\end{align*}
\begin{center}\small
\begin{tabular}{llrrrrrr}
\toprule
Case & Split & $L$ & $T$ & $D$ & $A$ & $\beta$ & Profile \\
\midrule
\texttt{heat-g0-v0} & train & 1 & 0.25 & 0.05 & 0.3 & 1 & sine \\
\texttt{heat-g0-v1} & train & 1.2 & 0.25 & 0.1 & 0.4 & 1 & sine \\
\texttt{heat-g0-v2} & train & 1 & 0.35 & 0.15 & 0.5 & 1 & sine \\
\texttt{heat-g0-v3} & train & 1.2 & 0.35 & 0.2 & 0.6 & 1 & sine \\
\texttt{heat-g1-v0} & validation & 1 & 0.25 & 0.05 & 0.3 & 2 & quadratic \\
\texttt{heat-g1-v1} & validation & 1.2 & 0.25 & 0.1 & 0.4 & 2 & quadratic \\
\texttt{heat-g1-v2} & validation & 1 & 0.35 & 0.15 & 0.5 & 2 & quadratic \\
\texttt{heat-g1-v3} & validation & 1.2 & 0.35 & 0.2 & 0.6 & 2 & quadratic \\
\bottomrule
\end{tabular}
\end{center}
\noindent Reserved cases: 4 (parameters kept evaluator-only).

\subsection{Convection--diffusion}
\begin{align*}
u_t + v\,u_x - D\,u_{xx} = f(x,t)\\
f = A\bigl[\beta t^{\beta-1}s - D\,g\,s'' + v\,g\,s'\bigr]
\end{align*}
\begin{center}\small
\begin{tabular}{llrrrrrrr}
\toprule
Case & Split & $L$ & $T$ & $D$ & $v$ & $A$ & $\beta$ & Profile \\
\midrule
\texttt{convection\_diffusion-g0-v0} & train & 1 & 0.25 & 0.05 & 0.25 & 0.3 & 1 & sine \\
\texttt{convection\_diffusion-g0-v1} & train & 1.2 & 0.25 & 0.1 & 0.35 & 0.4 & 1 & sine \\
\texttt{convection\_diffusion-g0-v2} & train & 1 & 0.35 & 0.15 & 0.45 & 0.5 & 1 & sine \\
\texttt{convection\_diffusion-g0-v3} & train & 1.2 & 0.35 & 0.2 & 0.55 & 0.6 & 1 & sine \\
\texttt{convection\_diffusion-g1-v0} & validation & 1 & 0.25 & 0.05 & 0.25 & 0.3 & 2 & quadratic \\
\texttt{convection\_diffusion-g1-v1} & validation & 1.2 & 0.25 & 0.1 & 0.35 & 0.4 & 2 & quadratic \\
\texttt{convection\_diffusion-g1-v2} & validation & 1 & 0.35 & 0.15 & 0.45 & 0.5 & 2 & quadratic \\
\texttt{convection\_diffusion-g1-v3} & validation & 1.2 & 0.35 & 0.2 & 0.55 & 0.6 & 2 & quadratic \\
\bottomrule
\end{tabular}
\end{center}
\noindent Reserved cases: 4 (parameters kept evaluator-only).

\subsection{Viscous Burgers}
\begin{align*}
u_t + u\,u_x - D\,u_{xx} = f(x,t)\\
f = A\bigl[\beta t^{\beta-1}s - D\,g\,s''\bigr] + A^2 g^2\,s\,s'
\end{align*}
\begin{center}\small
\begin{tabular}{llrrrrrr}
\toprule
Case & Split & $L$ & $T$ & $D$ & $A$ & $\beta$ & Profile \\
\midrule
\texttt{burgers-g0-v0} & train & 1 & 0.25 & 0.05 & 0.3 & 1 & sine \\
\texttt{burgers-g0-v1} & train & 1.2 & 0.25 & 0.1 & 0.4 & 1 & sine \\
\texttt{burgers-g0-v2} & train & 1 & 0.35 & 0.15 & 0.5 & 1 & sine \\
\texttt{burgers-g0-v3} & train & 1.2 & 0.35 & 0.2 & 0.6 & 1 & sine \\
\texttt{burgers-g1-v0} & validation & 1 & 0.25 & 0.05 & 0.3 & 2 & quadratic \\
\texttt{burgers-g1-v1} & validation & 1.2 & 0.25 & 0.1 & 0.4 & 2 & quadratic \\
\texttt{burgers-g1-v2} & validation & 1 & 0.35 & 0.15 & 0.5 & 2 & quadratic \\
\texttt{burgers-g1-v3} & validation & 1.2 & 0.35 & 0.2 & 0.6 & 2 & quadratic \\
\bottomrule
\end{tabular}
\end{center}
\noindent Reserved cases: 4 (parameters kept evaluator-only).

\subsection{Caputo subdiffusion}
\begin{align*}
\Caputo{\alpha}u - D\,u_{xx} = f(x,t),\quad 0<\alpha<1\\
f = A\Bigl[\frac{\Gamma(\beta+1)}{\Gamma(\beta+1-\alpha)}t^{\beta-\alpha}s - D\,g\,s''\Bigr]
\end{align*}
\begin{center}\small
\begin{tabular}{llrrrrrrr}
\toprule
Case & Split & $L$ & $T$ & $D$ & $A$ & $\beta$ & $\alpha$ & Profile \\
\midrule
\texttt{caputo-g0-v0} & train & 1 & 0.25 & 0.05 & 0.3 & 2 & 0.4 & sine \\
\texttt{caputo-g0-v1} & train & 1.2 & 0.25 & 0.1 & 0.4 & 2 & 0.55 & sine \\
\texttt{caputo-g0-v2} & train & 1 & 0.35 & 0.15 & 0.5 & 2 & 0.7 & sine \\
\texttt{caputo-g0-v3} & train & 1.2 & 0.35 & 0.2 & 0.6 & 2 & 0.85 & sine \\
\texttt{caputo-g1-v0} & validation & 1 & 0.25 & 0.05 & 0.3 & 0.8 & 0.4 & quadratic \\
\texttt{caputo-g1-v1} & validation & 1.2 & 0.25 & 0.1 & 0.4 & 0.8 & 0.55 & quadratic \\
\texttt{caputo-g1-v2} & validation & 1 & 0.35 & 0.15 & 0.5 & 0.8 & 0.7 & quadratic \\
\texttt{caputo-g1-v3} & validation & 1.2 & 0.35 & 0.2 & 0.6 & 0.8 & 0.85 & quadratic \\
\bottomrule
\end{tabular}
\end{center}
\noindent Reserved cases: 4 (parameters kept evaluator-only).

\subsection{Multi-term Caputo diffusion}
\begin{align*}
\sum_{j=1}^{2} w_j\,\Caputo{\alpha_j}u - D\,u_{xx} = f(x,t)\\
f = A\Bigl[\sum_j w_j\frac{\Gamma(\beta+1)}{\Gamma(\beta+1-\alpha_j)}t^{\beta-\alpha_j}s - D\,g\,s''\Bigr]
\end{align*}
\begin{center}\small
\begin{tabular}{llrrrrrrrr}
\toprule
Case & Split & $L$ & $T$ & $D$ & $A$ & $\beta$ & $(\alpha_1,\alpha_2)$ & $(w_1,w_2)$ & Profile \\
\midrule
\texttt{multi\_term-g0-v0} & train & 1 & 0.25 & 0.05 & 0.3 & 2 & (0.3, 0.7) & (0.4, 1) & sine \\
\texttt{multi\_term-g0-v1} & train & 1.2 & 0.25 & 0.1 & 0.4 & 2 & (0.35, 0.75) & (0.4, 1) & sine \\
\texttt{multi\_term-g0-v2} & train & 1 & 0.35 & 0.15 & 0.5 & 2 & (0.4, 0.8) & (0.4, 1) & sine \\
\texttt{multi\_term-g0-v3} & train & 1.2 & 0.35 & 0.2 & 0.6 & 2 & (0.45, 0.85) & (0.4, 1) & sine \\
\texttt{multi\_term-g1-v0} & validation & 1 & 0.25 & 0.05 & 0.3 & 0.8 & (0.3, 0.7) & (0.4, 1) & quadratic \\
\texttt{multi\_term-g1-v1} & validation & 1.2 & 0.25 & 0.1 & 0.4 & 0.8 & (0.35, 0.75) & (0.4, 1) & quadratic \\
\texttt{multi\_term-g1-v2} & validation & 1 & 0.35 & 0.15 & 0.5 & 0.8 & (0.4, 0.8) & (0.4, 1) & quadratic \\
\texttt{multi\_term-g1-v3} & validation & 1.2 & 0.35 & 0.2 & 0.6 & 0.8 & (0.45, 0.85) & (0.4, 1) & quadratic \\
\bottomrule
\end{tabular}
\end{center}
\noindent Reserved cases: 4 (parameters kept evaluator-only).


\subsection{Two-dimensional wave equation (8 public, 4 reserved cases)}
\begin{align*}
&u_{tt} - D\,(u_{xx}+u_{yy}) = f(x,y,t),\qquad (x,y)\in(0,L)^2,\ t\in(0,T],\\
&u = 0 \text{ on } \partial(0,L)^2,\qquad u(x,y,0)=A\,s(x)s(y),\qquad u_t(x,y,0)=0,\\
&u(x,y,t) = A\,s(x)\,s(y)\,\bigl(\cos\omega t + t^{\beta}\bigr),\qquad \omega=\sqrt{D}\,\pi\sqrt{2}/L,\\
&f = A\Bigl[\bigl(-\omega^2\cos\omega t+\beta(\beta-1)t^{\beta-2}\bigr)S
   - D\bigl(\cos\omega t+t^{\beta}\bigr)\Delta S\Bigr],\\
&S(x,y)=s(x)\,s(y),
\end{align*}
with $s$ as in the 1D benchmark (train: sine, $\beta=2.5$; validation: quadratic, $\beta=1.5$, so
$u_{tt}\sim t^{-1/2}$ is singular at $t=0$; reserved: mixed, $\beta=1.25$). $D=c^2$ is the squared wave speed.
\noindent Frequencies $\omega$: \texttt{wave2d-g0-v0}: 2.2214, \texttt{wave2d-g0-v1}: 2.7768, \texttt{wave2d-g0-v2}: 4.4429, \texttt{wave2d-g0-v3}: 4.6280, \texttt{wave2d-g1-v0}: 2.2214, \texttt{wave2d-g1-v1}: 2.7768, \texttt{wave2d-g1-v2}: 4.4429, \texttt{wave2d-g1-v3}: 4.6280.

\begin{center}\small
\begin{tabular}{llrrrrr}
\toprule
Case & Split & $L$ & $T$ & $D=c^2$ & $A$ & $\beta$ \\
\midrule
\texttt{wave2d-g0-v0} & train & 1 & 0.5 & 0.25 & 0.3 & 2.5 \\
\texttt{wave2d-g0-v1} & train & 1.2 & 0.5 & 0.5625 & 0.4 & 2.5 \\
\texttt{wave2d-g0-v2} & train & 1 & 0.75 & 1 & 0.5 & 2.5 \\
\texttt{wave2d-g0-v3} & train & 1.2 & 0.75 & 1.562 & 0.6 & 2.5 \\
\texttt{wave2d-g1-v0} & validation & 1 & 0.5 & 0.25 & 0.3 & 1.5 \\
\texttt{wave2d-g1-v1} & validation & 1.2 & 0.5 & 0.5625 & 0.4 & 1.5 \\
\texttt{wave2d-g1-v2} & validation & 1 & 0.75 & 1 & 0.5 & 1.5 \\
\texttt{wave2d-g1-v3} & validation & 1.2 & 0.75 & 1.562 & 0.6 & 1.5 \\
\bottomrule
\end{tabular}
\end{center}
\noindent Reserved cases: 4 (parameters kept evaluator-only).


\subsection{Group FPDE examples (16 public, 4 reserved cases)}
All start from rest and use $t\in(0,1]$. The fractional PDEs use $x\in(0,\pi)$ with $u(0,t)=u(\pi,t)=0$ and
$u(x,0)=0$. The exact solutions are $u=t^{p}$ (FODE) and $u=t^{p}\sin x$ (FPDEs). The paper instances are
$p=3$ with order coefficient $0.5$ (so $u=t^3$ or $u=t^3\sin x$); the other cases vary the order and $p$.

\subsection{Fractional ODE}
\begin{equation*}
\Caputo{\alpha}u(t) = f(t),\quad u(0)=0,\qquad u(t)=t^{p},\qquad
f(t)=\frac{\Gamma(p+1)}{\Gamma(p+1-\alpha)}\,t^{p-\alpha}.
\end{equation*}

\begin{center}\small
\begin{tabular}{llrr}
\toprule
Case & Split & $\alpha$ & $p$ \\
\midrule
\texttt{fode-g0-v0} & train & 0.3 & 2.5 \\
\texttt{fode-g0-v1} & train & 0.7 & 3.5 \\
\texttt{fode-g1-v0} & validation & 0.5 & 3 \\
\texttt{fode-g1-v1} & validation & 0.9 & 4 \\
\bottomrule
\end{tabular}
\end{center}
\noindent Reserved cases: 1 (parameters kept evaluator-only).

\subsection{Caputo FPDE}
\begin{equation*}
\Caputo{\alpha}u = u_{xx} + f(x,t),\qquad u=t^{p}\sin x,\qquad
f(x,t)=\Bigl[\frac{\Gamma(p+1)}{\Gamma(p+1-\alpha)}\,t^{p-\alpha} + t^{p}\Bigr]\sin x.
\end{equation*}

\begin{center}\small
\begin{tabular}{llrr}
\toprule
Case & Split & $\alpha$ & $p$ \\
\midrule
\texttt{fpde\_caputo-g0-v0} & train & 0.3 & 2.5 \\
\texttt{fpde\_caputo-g0-v1} & train & 0.7 & 3.5 \\
\texttt{fpde\_caputo-g1-v0} & validation & 0.5 & 3 \\
\texttt{fpde\_caputo-g1-v1} & validation & 0.9 & 4 \\
\bottomrule
\end{tabular}
\end{center}
\noindent Reserved cases: 1 (parameters kept evaluator-only).

\subsection{Time-variable-order FPDE}
\begin{align*}
&\Caputo{\alpha(t)}u = u_{xx} + f(x,t),\qquad \alpha(t)=c\,\sin^2 t,\qquad u=t^{p}\sin x,\\
&f=\Bigl[\frac{\Gamma(p+1)}{\Gamma(p+1-\alpha(t))}\,t^{p-\alpha(t)} + t^{p}\Bigr]\sin x.
\end{align*}

\begin{center}\small
\begin{tabular}{llrr}
\toprule
Case & Split & $c$ & $p$ \\
\midrule
\texttt{fpde\_variable\_order-g0-v0} & train & 0.3 & 2.5 \\
\texttt{fpde\_variable\_order-g0-v1} & train & 0.7 & 3.5 \\
\texttt{fpde\_variable\_order-g1-v0} & validation & 0.5 & 3 \\
\texttt{fpde\_variable\_order-g1-v1} & validation & 0.9 & 4 \\
\bottomrule
\end{tabular}
\end{center}
\noindent Reserved cases: 1 (parameters kept evaluator-only).

\subsection{Space--time-variable-order FPDE}
\begin{align*}
&\Caputo{\alpha(x,t)}u = u_{xx} + f(x,t),\qquad \alpha(x,t)=c\,t^{2}\sin x,\qquad u=t^{p}\sin x,\\
&f=\Bigl[\frac{\Gamma(p+1)}{\Gamma(p+1-\alpha(x,t))}\,t^{p-\alpha(x,t)} + t^{p}\Bigr]\sin x.
\end{align*}

\begin{center}\small
\begin{tabular}{llrr}
\toprule
Case & Split & $c$ & $p$ \\
\midrule
\texttt{fpde\_variable\_order\_xt-g0-v0} & train & 0.3 & 2.5 \\
\texttt{fpde\_variable\_order\_xt-g0-v1} & train & 0.7 & 3.5 \\
\texttt{fpde\_variable\_order\_xt-g1-v0} & validation & 0.5 & 3 \\
\texttt{fpde\_variable\_order\_xt-g1-v1} & validation & 0.9 & 4 \\
\bottomrule
\end{tabular}
\end{center}
\noindent Reserved cases: 1 (parameters kept evaluator-only).

\begin{thebibliography}{99}

\bibitem{Bas86}
C.~Basdevant, M.~Deville, P.~Haldenwang, J.~M.~Lacroix, J.~Ouazzani,
R.~Peyret, P.~Orlandi, and A.~T.~Patera,
Spectral and finite difference solutions of the Burgers equation,
\emph{Computers \& Fluids} 14(1) (1986) 23--41.

\bibitem{CH58}
J.~W.~Cahn and J.~E.~Hilliard,
Free energy of a nonuniform system. I. Interfacial free energy,
\emph{Journal of Chemical Physics} 28(2) (1958) 258--267.

\bibitem{Eva10}
L.~C.~Evans,
\emph{Partial Differential Equations}, 2nd ed.,
American Mathematical Society, Providence, RI, 2010.

\bibitem{Hab13}
R.~Haberman,
\emph{Applied Partial Differential Equations with Fourier Series and
Boundary Value Problems}, 5th ed., Pearson, Boston, 2013.

\bibitem{Jac99}
J.~D.~Jackson,
\emph{Classical Electrodynamics}, 3rd ed., Wiley, New York, 1999.

\bibitem{KT05}
A.-K.~Kassam and L.~N.~Trefethen,
Fourth-order time-stepping for stiff PDEs,
\emph{SIAM Journal on Scientific Computing} 26(4) (2005) 1214--1233.

\bibitem{Kov48}
L.~I.~G.~Kovasznay,
Laminar flow behind a two-dimensional grid,
\emph{Mathematical Proceedings of the Cambridge Philosophical Society}
44(1) (1948) 58--62.

\bibitem{Mit13}
W.~F.~Mitchell,
A collection of 2D elliptic problems for testing adaptive grid
refinement algorithms,
\emph{Applied Mathematics and Computation} 220 (2013) 350--364.

\bibitem{Mon03}
P.~Monk,
\emph{Finite Element Methods for Maxwell's Equations},
Oxford University Press, Oxford, 2003.

\bibitem{Pat80}
S.~V.~Patankar,
\emph{Numerical Heat Transfer and Fluid Flow},
Hemisphere, Washington, DC, 1980.

\bibitem{RPK19}
M.~Raissi, P.~Perdikaris, and G.~E.~Karniadakis,
Physics-informed neural networks: A deep learning framework for solving
forward and inverse problems involving nonlinear partial differential
equations,
\emph{Journal of Computational Physics} 378 (2019) 686--707.

\bibitem{RST08}
H.-G.~Roos, M.~Stynes, and L.~Tobiska,
\emph{Robust Numerical Methods for Singularly Perturbed Differential
Equations}, 2nd ed., Springer, Berlin, 2008.

\bibitem{Str08}
W.~A.~Strauss,
\emph{Partial Differential Equations: An Introduction}, 2nd ed.,
Wiley, Hoboken, NJ, 2008.

\bibitem{TG37}
G.~I.~Taylor and A.~E.~Green,
Mechanism of the production of small eddies from large ones,
\emph{Proceedings of the Royal Society of London A} 158 (1937) 499--521.

\bibitem{WWP21}
S.~Wang, H.~Wang, and P.~Perdikaris,
On the eigenvector bias of Fourier feature networks: From regression to
solving multi-scale PDEs with physics-informed neural networks,
\emph{Computer Methods in Applied Mechanics and Engineering} 384 (2021)
113938.

\bibitem{Whi74}
G.~B.~Whitham,
\emph{Linear and Nonlinear Waves}, Wiley, New York, 1974.

\bibitem{Yee66}
K.~S.~Yee,
Numerical solution of initial boundary value problems involving
Maxwell's equations in isotropic media,
\emph{IEEE Transactions on Antennas and Propagation} 14(3) (1966) 302--307.

\end{thebibliography}
\end{document}
